\chapter{Conclusiones}

La implementación de los algoritmos de \textit{Divide y Vencerás} permitió comparar su comportamiento real con la complejidad teórica. Los benchmarks muestran que esta estrategia es más eficiente que los enfoques iterativos, aunque su rendimiento también depende de decisiones como el pivote, el umbral y el tamaño de entrada.

En ordenamiento, \textit{Merge Sort} y \textit{Quick Sort} superaron con claridad a los algoritmos iterativos. \textit{Merge Sort} fue el más estable, y su versión optimizada mejoró ligeramente el rendimiento al usar \textit{Insertion Sort} en subarreglos pequeños.
En \textit{Quick Sort}, la elección del pivote fue decisiva. Los pivotes fijo primero y último elemento degradaron el rendimiento en los peores casos, mientras que las variantes aleatoria y mediana de tres mantuvieron mejor estabilidad. La mediana de tres resultó la opción más robusta.

Las búsquedas binaria iterativa, binaria recursiva y exponencial mostraron tiempos muy bajos y funcionaron mejor que la búsqueda secuencial. Además, la versión recursiva de búsqueda binaria no mostró una penalización importante frente a la iterativa.

En selección, \textit{Quick Select} se comportó de forma cercana a lo esperado para un algoritmo lineal. Como en \textit{Quick Sort}, hubo diferencia entre mejor y peor caso, aunque la estrategia de mediana de tres ayuda a reducir esa variación.

El benchmark de umbral mostró que existe un punto intermedio conveniente entre dividir demasiado y ordenar demasiado pronto con un algoritmo simple. El mejor resultado se obtuvo con un umbral de 32 elementos, por lo que ese valor quedó como una buena elección para esta implementación.

En general, los resultados confirman que \textit{Divide y Vencerás} mejora el rendimiento frente a enfoques más simples. Sin embargo, el comportamiento real también depende de detalles de implementación y de cómo están distribuidos los datos, así que siempre conviene medir antes de fijar una solución definitiva.

\chapter{Contribuciones}

\begin{itemize}
    \item Iván Mansilla: Integración general, búsqueda exponencial, búsqueda binaria por rango, benchmarking, informe.
    \item Diego Peralta: QuickSort, búsqueda binaria recursiva, N mejores deportistas, informe.
    \item Catalina Viñas: Merge Sort, búsqueda por interpolación, QuickSelect, gráficos.
\end{itemize}


